% ---------------------------------------------------------------------
%  Chapitre 3 — JUGER LES STRATÉGIES (le cœur du mémoire)
%  Suit les 5 points des consignes pour CHAQUE méthode :
%   théorie -> conditions d'application -> résultats -> limites
%  Structure : étape 1 (test de référence, rédigée) -> étape 1 bis
%  (l'escalade de prudence face à la dépendance) -> lecture du verdict
%  (le marché comme série temporelle) -> synthèse.
% ---------------------------------------------------------------------
\chapter{Juger les stratégies malgré la dépendance~: une escalade de prudence}
\label{chap:juger}
\label{chap:bloc1} % ancien nom du label conservé pour d'éventuels renvois

Ce chapitre traite une question simple à formuler et difficile à trancher.
Parmi les dix stratégies d'analyse technique du banc d'essai, auxquelles
s'ajoute le témoin, lesquelles produisent un avantage significativement
positif~? Le test de référence est d'abord appliqué sous ses hypothèses
habituelles, puis ses conditions d'application sont mises à l'épreuve une à
une, en s'appuyant sur la structure de dépendance mesurée au
chapitre~\ref{chap:donnees}.

% =====================================================================
\section{Étape 1 --- chaque stratégie bat-elle le hasard ?}
\label{sec:b1-etape1}

\subsection{La question posée à l'outil}

La première question posée à l'outil est de savoir si l'avantage moyen
affiché par une stratégie est réel, ou s'il n'est qu'un accident dû à
l'échantillon dont nous disposons. On la formalise en confrontant la moyenne
de l'\edgevar{} à la valeur zéro, ce qui conduit au test de Student présenté
ci-dessous.

\subsection{Les méthodes qui y répondent}

% --- MÉTHODE 1 : le test principal -----------------------------------
\subsubsection*{Test de Student}

Lorsqu'une stratégie affiche un \edgevar{} moyen positif, on cherche à savoir
si elle a réellement capté une information annonciatrice, ou si ce résultat
tient au hasard de l'échantillon.

Il faut ici distinguer deux quantités. L'\edgevar{} moyen théorique
\(\mu_{\edgevar}\) est la vraie valeur, inconnue, qui gouverne la stratégie,
tandis que la moyenne observée \(\bar e\) est celle que l'on calcule sur notre
échantillon d'opérations, et qui est connue. La question se pose alors ainsi.
À partir de la valeur observée \(\bar e\), peut-on conclure que la vraie
valeur \(\mu_{\edgevar}\) est positive~?

Le test de Student à un échantillon répond à cette question en confrontant
\(\mu_{\edgevar}\) à une valeur de référence fixée à zéro, zéro signifiant
ici l'absence de tout avantage sur le hasard. Sa mise en \oe uvre suppose deux
conditions, à savoir que l'\edgevar{} vérifie un critère de normalité et que
les opérations soient indépendantes les unes des autres. La première sera
assouplie par le théorème central limite présenté plus loin, la seconde, plus
délicate, sera examinée à la section suivante. On oppose ainsi deux
hypothèses,
\[
  \begin{cases}
    \Hzero : \mu_{\edgevar} = 0 & \text{(le vrai avantage est nul, pas mieux que le hasard),}\\[2pt]
    \Hun : \mu_{\edgevar} > 0 & \text{(le vrai avantage est positif).}
  \end{cases}
\]
Le test est unilatéral, car seul un avantage strictement positif nous
intéresse.

Pour trancher, on définit \(t\) comme le nombre d'erreurs-types nécessaires
pour faire passer l'\edgevar{} moyen de zéro à la valeur obtenue dans notre
échantillon,
\[
  t = \frac{\bar e}{s/\sqrt{n}} \sim t_{n-1} \quad\text{sous } \Hzero .
\]

Sous \Hzero{}, c'est-à-dire lorsque le vrai avantage est nul, \(t\) suit une
loi de Student à \(n-1\) degrés de liberté. On évalue alors la p-value, qui
est la probabilité d'obtenir un \(t\) au moins aussi grand que celui observé
si \Hzero{} était vraie.

Si cette p-value est très faible, le hasard cesse d'être une explication
crédible de l'avantage moyen constaté. On conclut alors que l'\edgevar{}
moyen théorique est positif, et donc que la stratégie fait mieux que le
hasard.

Il a été dit plus haut que le test de Student exige que l'\edgevar{} suive une
loi normale. Cette condition sera examinée en deux temps, d'abord par le test
de Shapiro et Wilk, puis à l'aide du théorème central limite. Le premier a peu
de chances de valider l'hypothèse, car la taille considérable de notre
échantillon le rend extrêmement sensible, alors que cette même taille renforce
au contraire le second.

% --- MÉTHODE 2 : la condition du test principal ----------------------
\subsubsection*{Test de normalité de Shapiro et Wilk}

Ce test de normalité compare deux mesures de dispersion, d'une part les
écarts à la moyenne, d'autre part l'ordre des valeurs. Ces deux mesures ne
coïncident que si l'hypothèse de normalité est vraie. On oppose ainsi deux
hypothèses,
\[
  \begin{cases}
    \Hzero \;(W \simeq 1) : \text{l'\edgevar{} suit une loi normale,}\\[2pt]
    \Hun \;(W \ll 1) : \text{l'\edgevar{} ne suit pas une loi normale.}
  \end{cases}
\]
Contrairement au test de Student, c'est ici \Hzero{} qui porte la propriété
recherchée, à savoir la normalité, si bien qu'une p-value faible conduira à
rejeter celle-ci.

On pose alors la statistique \(W = b^2 / \mathrm{SCE}\) comme rapport de ces
deux mesures. Lorsque \(W\) est proche de 1, l'échantillon est compatible avec
une loi normale.

Pour trancher, on procède comme pour le test de Student. On calcule la
probabilité d'obtenir notre \(W\) observé si la loi était normale, et si cette
probabilité est trop faible, l'hypothèse de normalité est rejetée.

Appliqué à nos données, le test rejette la normalité pour les onze
stratégies, comme le montre le tableau~\ref{tab:shapiro}. La condition du test
de Student n'est donc, en apparence, pas remplie. On remarque que même les
stratégies dont le \(W\) est proche de 1, comme \texttt{dt\_top} avec
$W = 0{,}95$, sont rejetées. À grand effectif, le moindre écart à la normale
suffit à produire une p-value négligeable, ce qui justifie de recourir au
théorème central limite plutôt qu'à ce test pour valider l'emploi de Student.

\begin{table}[H]
\centering
\caption{Statistique \(W\) et p-value de Shapiro--Wilk par stratégie
(\edgevar{}~; sous-échantillon de 5\,000 opérations lorsque \(n > 5\,000\)).}
\label{tab:shapiro}
\small
\begin{tabular}{@{}lrrr@{}}
\toprule
Stratégie & $n$ & $W$ & $p$ (Shapiro) \\
\midrule
\texttt{oracle} (témoin) & 6\,798 & $0{,}590$ & $6{,}8\times10^{-76}$ \\
\midrule
\texttt{db\_bottom}    & 8\,513 & $0{,}911$ & $2{,}7\times10^{-47}$ \\
\texttt{dt\_top}       & 7\,381 & $0{,}950$ & $3{,}2\times10^{-38}$ \\
\texttt{hs\_classic}   & 1\,065 & $0{,}886$ & $2{,}4\times10^{-27}$ \\
\texttt{hs\_inverse}   & 1\,206 & $0{,}963$ & $6{,}4\times10^{-17}$ \\
\texttt{ma\_crossover} & 5\,417 & $0{,}498$ & $3{,}2\times10^{-80}$ \\
\texttt{rsi\_classic}  & 7\,117 & $0{,}864$ & $1{,}6\times10^{-54}$ \\
\texttt{rsi\_strict}   &   978  & $0{,}717$ & $1{,}1\times10^{-37}$ \\
\texttt{rsi\_trend}    & 2\,789 & $0{,}883$ & $1{,}2\times10^{-41}$ \\
\texttt{sr\_breakdown} & 1\,377 & $0{,}775$ & $8{,}0\times10^{-40}$ \\
\texttt{sr\_breakout}  & 2\,367 & $0{,}843$ & $2{,}1\times10^{-43}$ \\
\bottomrule
\end{tabular}
\end{table}

% --- LE TCL : ce qui débloque le test de Student ---------------------
\subsubsection*{Théorème central limite}

Le rejet prononcé par le test précédent porte sur la loi de l'\edgevar{} pris
individuellement, opération par opération. Or le test de Student ne dépend
que de la moyenne de ces \edgevar{}, et non de chacun d'eux. On peut donc se
rabattre sur le théorème central limite, qui garantit la quasi-normalité de
la moyenne sur les grands échantillons, quelle que soit la loi de départ,
\[
  \frac{\bar e - \mu_{\edgevar}}{\sigma/\sqrt{n}}
  \;\xrightarrow[n \to \infty]{}\; \mathcal{N}(0,1) .
\]
Il faut insister sur ce point, car la nuance est décisive. Le théorème ne dit
pas que les données sont normales, il dit que leur moyenne l'est. C'est
précisément cette normalité de la moyenne dont le test de Student a besoin.

Cette justification suppose évidemment un effectif suffisant, et un outil
générique devrait écarter les stratégies trop peu fournies. Ici, \(n\) est
compris entre 978 et 8\,513 selon la stratégie, ce qui est largement
suffisant. La condition réellement utile au test de Student est donc
satisfaite, malgré le rejet du test de normalité.

On remarque de plus que la quantité ci-dessus est exactement celle du test de
Student. Sous \Hzero{}, où \(\mu_{\edgevar} = 0\), et en remplaçant
l'écart-type théorique \(\sigma\), inconnu, par son estimation \(s\), on
retrouve la statistique \(t = \bar e / (s/\sqrt{n})\). Le théorème central
limite justifie ainsi directement la forme et la validité du test.

% --- MÉTHODE 3 : le garde-fou des tests multiples --------------------
\subsubsection*{Correction de Bonferroni}

Pour appliquer le test de Student, il faut fixer un seuil au-dessous duquel
la p-value entraînera le rejet. Lorsqu'on mène un test isolé, on retient
habituellement \(0{,}05\), considérant que s'il n'y avait que cinq chances sur
cent d'obtenir un tel résultat sous \Hzero{}, il devient peu vraisemblable d'y
être tombé par hasard. On accepte alors le risque de rejeter \Hzero{} à tort.

Seulement, nous menons ici onze tests. Plaçons-nous dans le cas le plus
défavorable, celui où aucune stratégie ne possède d'avantage réel, \Hzero{}
étant vraie pour les onze. Pour un test isolé, la probabilité de ne pas
conclure à tort à un avantage vaut alors \(0{,}95\). En supposant les tests
indépendants, la probabilité de ne commettre aucune erreur sur l'ensemble
tombe à \(0{,}95^{11} \approx 0{,}57\). Le risque de commettre au moins un
faux positif s'élève donc à \(1 - 0{,}95^{11} \approx 0{,}43\), ce qui est
inacceptable.

Pour l'éviter, on raisonne en sens inverse, en fixant le risque global à
\(0{,}05\) et en déduisant le seuil de chaque test. On s'appuie pour cela sur
l'inégalité de Boole, selon laquelle la probabilité qu'au moins un événement
se produise est inférieure ou égale à la somme de leurs probabilités
individuelles. Pour garantir un risque global inférieur à \(0{,}05\), il
suffit donc que la somme des seuils individuels vaille \(0{,}05\). En les
prenant tous égaux, on obtient la correction de Bonferroni,
\[
  \alpha' = \frac{\alpha}{m} = \frac{0{,}05}{11} \approx 0{,}0045 .
\]

\subsection{Résultats}
% Chiffres RÉELS extraits de bloc1/03_resultats/etape1_tests.txt
\begin{table}[H]
\centering
\caption{Test de l'avantage moyen par stratégie (seuil Bonferroni $\approx 0{,}0045$).}
\label{tab:b1-etape1}
\small
\begin{tabular}{@{}lrrrc@{}}
\toprule
Stratégie & $n$ & \edgevar{} moyen & $p$ (Student) & Verdict \\
\midrule
\texttt{oracle} (témoin) & 6\,798 & $+0{,}0658$ & $\approx 0$    & \textbf{bat le hasard} \\
\midrule
\texttt{rsi\_classic} & 7\,117 & $+0{,}0123$ & $1{,}2\times10^{-8}$  & \textbf{bat le hasard} \\
\texttt{rsi\_strict}  &   978  & $+0{,}1313$ & $5{,}5\times10^{-10}$ & \textbf{bat le hasard} \\
\texttt{rsi\_trend}   & 2\,789 & $+0{,}0021$ & $0{,}25$             & non \\
\texttt{db\_bottom}   & 8\,513 & $-0{,}0025$ & $0{,}98$             & non \\
\texttt{ma\_crossover}& 5\,417 & $-0{,}0248$ & $\approx 1$          & non \\
\multicolumn{5}{c}{\footnotesize (autres stratégies~: $p>0{,}45$, non significatives)} \\
\bottomrule
\end{tabular}
\end{table}

Outre l'oracle témoin, qui bat le hasard de très loin comme on pouvait s'y
attendre, seules \texttt{rsi\_classic} et \texttt{rsi\_strict} survivent à la
correction de Bonferroni.

Il est utile de rappeler ce que ce test établit exactement. Il répond
uniquement à la question de savoir si l'\edgevar{} moyen positif s'explique
par le seul échantillon, auquel cas on conserve \Hzero{} et la stratégie n'est
pas retenue. Dans le cas contraire, on pourra affirmer que la stratégie a
capté un signal, mais non qu'elle est gagnante, et ce pour deux raisons.

D'une part, un \edgevar{} positif ne signifie pas des gains positifs, mais
seulement que l'on perd moins que le hasard. D'autre part, le grand nombre
d'observations écrase la largeur de la loi normale de la moyenne. On fait
alors ressortir comme significatives des stratégies dont l'\edgevar{} moyen
n'est que très faiblement positif, et que les frais de transaction
suffiraient à rendre perdantes.




% =====================================================================
\section{Étape 1 bis --- le verdict survit-il à la dépendance ?}
\label{sec:b1-dependance}

\subsection{La question posée à l'outil}

Le test de Student a été appliqué de façon naïve, en supposant les
observations indépendantes. Or, comme nous l'avons mesuré au
chapitre~\ref{chap:donnees}, nos données sont dépendantes par le titre, par le
marché et d'un mois à l'autre.

Ces dépendances n'empruntent en réalité que deux chemins, que nous appellerons
les deux portes. La première est la porte de l'action, celle des opérations
ouvertes sur un même titre. La seconde est la porte de la période, celle des
opérations menées à une même époque, qu'elles subissent ensemble le mouvement
du marché ou qu'elles se prolongent sur les mêmes mois. Ce sont les deux
seules voies par lesquelles deux opérations peuvent partager une cause
commune.

Le problème est que le test de Student traite les opérations d'une stratégie
comme autant d'informations distinctes, alors que ces informations se
répètent en partie. L'erreur-type s'en trouve sous-estimée. Comme elle sert à
mesurer l'écart entre l'avantage observé et la valeur de référence, un
dénominateur trop petit gonfle la statistique et rend la p-value trop
optimiste. Le risque de conclure à tort qu'une stratégie bat le hasard
augmente d'autant.

Nous cherchons donc à savoir si le verdict positif obtenu pour
\texttt{rsi\_classic} et \texttt{rsi\_strict} résiste lorsque l'on abandonne
l'hypothèse d'indépendance. La démarche suit quatre temps. Nous corrigeons
d'abord l'erreur-type au moyen de l'effet de sondage, puis nous vérifions ce
résultat par un rééchantillonnage qui ne suppose aucune forme particulière des
données. Nous fermons ensuite chaque porte en regroupant les opérations, de
sorte que chaque titre, puis chaque période, ne compte qu'une fois. Nous
vérifions enfin que la porte de la période tient elle-même, en traitant la
suite des avantages mensuels comme une série chronologique.

Une règle de décision gouverne cette escalade, et il faut l'énoncer avant de
la parcourir. Chaque étage traite une faille laissée ouverte par le
précédent, et l'échec à un seul d'entre eux suffit à écarter une stratégie de
façon définitive. Le succès, en revanche, n'accorde jamais de blanc-seing, il
ouvre seulement le droit de passer à l'étage suivant. Le témoin, lui, doit les
franchir tous, faute de quoi ce serait l'outil qui serait trop sévère, et non
les stratégies qui seraient mauvaises.

\subsection{Les méthodes qui y répondent}

% --- MÉTHODE 1 : corriger la variance (cours de sondages) -------------
\subsubsection*{Corriger l'erreur-type par l'effet de sondage}

Comme nous l'avons vu au chapitre précédent, les titres forment des grappes
dans nos données. Nous avons également mesuré la corrélation intra-titre
$\hat\rho$, qui dit à quel point les opérations portant sur un même titre se
ressemblent.

Introduisons maintenant l'effet de sondage (\emph{design effect}), qui est le
rapport entre la variance d'un plan de sondage et celle d'un plan de
référence, ce dernier étant le sondage aléatoire simple, où chaque individu
est tiré indépendamment des autres. Ici, ce sera donc le rapport entre la
variance réelle, celle qui subit le regroupement des opérations par titre, et
la variance qu'on obtiendrait sous indépendance,
\[
  \mathrm{DEFF}
    = \frac{\mathrm{Var}(\hat\theta)}{\mathrm{Var}_{\mathrm{SAS}}(\hat\theta)} .
\]

Notre situation est celle d'un sondage à plusieurs degrés, et le cours de
théorie des sondages établit qu'alors l'effet de sondage vaut
\[
  \mathrm{DEFF} = 1 + (\bar m - 1)\,\hat\rho ,
\]
où $\bar m$ désigne le nombre moyen d'opérations par titre et $\hat\rho$ la
corrélation intra-titre mesurée au chapitre~\ref{chap:donnees}. La lecture de
cette formule est éclairante. Si $\hat\rho$ est nul, l'effet de sondage vaut
un et rien ne change. S'il est positif, l'effet dépasse un, et d'autant plus
que les grappes sont grosses, car entasser des opérations qui se ressemblent
dans un même titre aggrave la redondance.

La variance réelle de la moyenne vaut donc $\mathrm{DEFF}$ fois celle que l'on
obtiendrait sous indépendance. Tout se passe comme si nous ne disposions que
de $n_{\mathrm{eff}}$ informations réellement distinctes,
\[
  n_{\mathrm{eff}} = \frac{n}{\mathrm{DEFF}} .
\]

Les largeurs d'intervalle de confiance se trouvent alors multipliées par
$\sqrt{\mathrm{DEFF}}$ par rapport au plan de référence, ce qui donne
l'erreur-type corrigée
\[
  \widehat{\mathrm{se}}_{\mathrm{corr}}
        = \frac{s}{\sqrt{n}}\sqrt{\mathrm{DEFF}} .
\]

On reprend enfin le test de Student de l'étape 1 avec cette erreur-type
corrigée, en conservant le même seuil de Bonferroni. Une seule quantité change
donc, le dénominateur, mais il suffit à rendre le test plus exigeant.

% --- MÉTHODE 2 : robustesse sans hypothèse de forme -------------------
\subsubsection*{Vérifier sans hypothèse de forme}

Le test précédent s'appuie encore sur la loi de Student, donc sur une forme
supposée des données. La méthode qui suit s'en affranchit. Elle ne cherche pas
à tester la dépendance intra-titre, mais à évaluer son effet sur l'\edgevar{}
moyen.

Plutôt que d'invoquer une loi théorique, on construit directement, à partir
des grappes, la courbe de répartition de l'\edgevar{} moyen. L'idée est la
suivante. Si la dépendance pèse lourd, cette courbe sera large et la
probabilité de tomber sous zéro augmentera d'autant.

Pour produire cette courbe, on définit comme précédemment chaque titre comme
une grappe à l'intérieur de laquelle les données sont dépendantes. On effectue
alors un grand nombre de tirages, ici 4\,000. À chaque tirage, on prélève avec
remise autant de titres que la stratégie en a mobilisés, soit environ cinq
cents. Un même titre peut sortir plusieurs fois, auquel cas on prend autant de
fois l'ensemble des opérations qu'il contient. On calcule ensuite la moyenne
des \edgevar{} de toutes les opérations ainsi rassemblées.

Le point décisif tient à cette façon de tirer. En prélevant des titres
entiers plutôt que des opérations isolées, on transporte intacte la structure
de dépendance interne de chaque grappe. Tirer opération par opération
mélangerait les titres et détruirait précisément ce que l'on cherche à
mesurer.

Après les 4\,000 tirages, on obtient une courbe dont la moyenne ne diffère pas
de la moyenne initiale. Le tirage étant uniforme sur les titres, chaque
opération a la même chance d'être reprise, si bien que la moyenne des tirages
se recentre sur celle de départ. Ce n'est donc pas le centre qui nous
intéresse, mais la largeur, car c'est elle qui traduit l'effet de la
dépendance intra-titre sur la fiabilité de notre moyenne.

Il faut préciser que l'on calcule la moyenne sur l'ensemble des opérations
obtenues, ce qui conserve à chaque titre le poids que lui donne son nombre
d'opérations, contrairement à la méthode de la porte présentée plus loin.
C'est justement en gardant ce poids que l'on conserve l'information portée par
la dépendance intra-titre.

Lorsque la dépendance pèse lourd, la courbe s'élargit et déborde sous zéro. On
compte alors le nombre de tirages passant sous zéro, que l'on divise par le
nombre total de tirages. Cette proportion tient lieu de p-value, et si elle
dépasse le seuil défini par la correction de Bonferroni, la stratégie est
rejetée.

Cette technique de rééchantillonnage est le \emph{bootstrap}, proposé par
Efron~\cite{efron1979}, dans sa variante par grappes qui préserve la
dépendance interne des blocs.

Elle suppose néanmoins que les titres soient indépendants entre eux, ce qui
n'est pas le cas ici en raison du facteur marché. Elle reste toutefois un
premier filtre légitime. En ignorant cette dépendance, elle sous-estime la
largeur de la courbe et se montre donc trop permissive, si bien qu'une
stratégie échouant déjà à ce stade est écartée sans appel, tandis que les
autres restent à départager.

Une réserve importante subsiste pourtant. En conservant le poids de chaque
titre, la méthode nous laisse vulnérables. Si la grande majorité des
opérations provient d'un petit nombre de titres, ces titres seront présents
dans presque tous les tirages, qui se ressembleront alors tous, et la courbe
sera étroite pour une mauvaise raison. Une courbe étroite ne suffit donc pas à
conclure que l'avantage est réel. Pour lever ce doute, il faut retirer aux
gros titres le poids que leur donne leur nombre d'opérations, ce qui est
l'objet de la méthode suivante.

% --- MÉTHODE 3 : le test décisif -------------------------------------
\subsubsection*{Fermer chaque porte en donnant un vote à chaque unité}

Nous avons mesuré trois canaux de dépendance, qui empruntent deux portes. La
première est celle de l'action, la seconde celle de la période, laquelle
englobe à la fois la dépendance de marché et le chevauchement des opérations.

On emploie ici la méthode de l'agrégation, qui consiste à calculer l'avantage
moyen de chaque grappe, puis à appliquer le test de Student à ces moyennes,
chaque grappe ne disposant que d'une seule voix quel que soit le nombre
d'opérations qu'elle contient. Pour la porte de l'action, on agrège par titre.
Pour la porte de la période, on agrège par mois, chaque opération étant
rattachée à son mois d'entrée, celui où le signal s'est déclenché et où la
décision a été prise. On évite ainsi que la fluctuation du marché sur une même
période, ou le chevauchement des positions, ne soient comptés comme autant
d'informations distinctes.

Dans les deux cas, la dépendance à l'intérieur d'une grappe ne peut plus rien
fausser, puisqu'un titre, ou un mois, ne pèse qu'une seule voix. En revanche,
rattacher une opération à son seul mois d'entrée alors qu'elle peut rester
ouverte plusieurs mois lie les mois voisins entre eux, et cette dépendance
d'un mois à l'autre fera l'objet de la méthode suivante.

Le théorème central limite s'applique encore à ces moyennes de grappes, dont
le nombre reste élevé. Selon les stratégies, il subsiste de 431 à 501~titres
pour la porte de l'action, et de 138 à 197~mois pour celle de la période, ce
qui est amplement suffisant. Le protocole ayant vocation à juger les
stratégies à venir, il signale toutefois le cas où une stratégie trop peu
diffusée laisserait un nombre de grappes insuffisant pour invoquer le
théorème. Le verdict serait alors tenu pour ininterprétable, et non pour
défavorable. Pour être validée, une stratégie doit franchir les deux portes,
chacune au seuil de Bonferroni.

Cette technique fait perdre beaucoup d'information, et il faut l'assumer. Nous
privilégions le risque de rejeter une bonne stratégie à celui d'en retenir une
mauvaise, ce qui est le choix propre à un outil de validation.

La dépendance intra-titre ayant déjà été traitée par l'effet de sondage et le
rééchantillonnage, l'apport propre de cette méthode est ailleurs. Elle met en
lumière les stratégies qui, comme nous l'avons expliqué à la fin de la méthode
précédente, ne reposent que sur quelques titres, car l'agrégation supprime le
poids que leur donnait leur nombre d'opérations.

Les stratégies qui ne franchissent pas les portes sont écartées. Une stratégie
rejetée parce que son avantage tient à quelques titres particuliers pourra
toutefois être réétudiée, afin d'affiner la sélection des opérations. Il n'est
pas exclu qu'une stratégie ne soit efficace que sur un secteur donné. Une
telle réétude devra évidemment se garder du sur-apprentissage, qui consiste à
repérer après coup, sur les données mêmes, une régularité que le hasard
suffisait à produire.

Une objection mérite d'être levée. La porte de l'action suppose les titres
indépendants entre eux, ce qui est contestable puisque le marché explique une
grande part de la variance et que deux titres travaillés au même moment
subissent le même choc. Mais cette dépendance est de nature temporelle, deux
titres n'étant liés que s'ils sont travaillés à la même période, et c'est
précisément ce dont la porte de la période se charge. Les deux portes se
complètent donc, chacune prenant en charge une source, ce qui justifie de les
poser toutes les deux.

Elles ne sont pourtant pas fermées simultanément. Nous conduisons deux tests
séparés plutôt qu'une agrégation croisée par titre et par mois, laquelle
produirait des grappes bien plus petites et rouvrirait à moitié chacune des
deux portes. L'exigence conjointe est d'ailleurs plus sévère, puisqu'une
stratégie doit survivre à chaque source de dépendance prise isolément et
traitée dans le pire des cas. Seule échapperait une dépendance propre à
l'intersection, un effet qui n'existerait que pour un titre donné pendant un
mois donné, et l'on voit mal comment un tel effet porterait à lui seul un
avantage. Le traitement simultané des deux dimensions existe, sous le nom de
double regroupement (\emph{double clustering}) de Petersen~\cite{petersen2009},
mais il sort du cadre des méthodes simples retenues ici. Nous y revenons en
section~\ref{sec:limites}.

Ce principe d'une statistique par grappe, puis d'une inférence sur la suite de
ces statistiques, appliqué à la porte de la période, est exactement la démarche
de Fama et MacBeth~\cite{famamacbeth1973}, standard de l'économétrie financière.



% --- MÉTHODE 4 : la porte « mois » est-elle vraiment fermée ? ---------
\subsubsection*{Vérifier que la porte de la période tient elle-même}

Comme nous l'avons mesuré au deuxième canal du chapitre~\ref{chap:donnees}, le
marché explique 73~\% de la variance, et lorsqu'un événement survient, tout le
marché réagit de façon semblable. L'agrégation mensuelle regroupe ainsi les
opérations déclenchées par un même contexte, et rattacher chacune à son seul
mois d'entrée garantit qu'aucune ne pèse dans deux grappes à la fois.

Ce rattachement règle le double comptage, mais non le recouvrement des
périodes de détention. Nous avons vu que les opérations se prolongent souvent
au-delà d'un mois. Une position ouverte en janvier et close en avril est rangée
dans janvier, alors que son résultat dépend de ce qui s'est passé jusqu'en
avril. La grappe de janvier porte donc de l'information sur les mois suivants,
information que la grappe de février porte également. Aucune opération n'est
comptée deux fois, et pourtant deux mois voisins partagent les mêmes conditions
de marché. Avec 73~\% de la variance expliquée par ces conditions communes, ce
recouvrement suffit à les faire se ressembler.

La suite des avantages mensuels est donc autocorrélée, et le test de Student
appliqué aux moyennes mensuelles pourrait rester trop optimiste. On voudrait
alors reprendre la méthode de l'effet de sondage pour corriger l'erreur-type.
Mais nos données sont ici chronologiques, et la dépendance y est ordonnée, un
mois étant davantage lié au précédent qu'à celui d'il y a cinq mois, là où les
opérations d'un même titre étaient interchangeables. Un coefficient unique ne
peut donc pas la décrire, et il nous faut une autre méthode.

Nous voulons calculer, pour un mois~$t$, l'impact moyen des mois qui le
précèdent, en remontant jusqu'à six mois en arrière. Pour cela, la suite des
avantages mensuels sera traitée comme une série chronologique. On supposera
qu'aucun effet de saison n'intervient, et surtout que le lien entre deux mois
ne dépend que de leur écart et non de leur date, le mois de janvier se
comportant comme les autres mois de l'année.

On mesure alors l'impact propre de chaque décalage par l'autocorrélation
partielle, qui retranche l'influence des mois intermédiaires. Un mois peut en
effet sembler lié à celui d'il y a deux mois par le seul fait qu'ils sont tous
deux liés au mois qui les sépare. Cette autocorrélation partielle s'annule
au-delà du dernier mois qui compte vraiment, ce qui nous donne le nombre de
mois à retenir, puis le poids de chacun. On corrigera ensuite l'erreur-type du
test mensuel, comme nous l'avions fait avec l'effet de sondage.

La méthode se déroule en trois temps, chacun porté par un outil. On établit
d'abord s'il existe réellement une dépendance, par le test de Box et Pierce.
On examine ensuite si l'on peut la modéliser, ce qui suppose la stationnarité,
vérifiée par les tests de Dickey et Fuller puis de Kwiatkowski et ses
coauteurs, l'ordre étant donné par l'autocorrélation partielle et les poids
par les équations de Yule et Walker. On détermine enfin de combien corriger,
au moyen de la variance de long terme. Les tests de conditions qui jalonnent
ce parcours sont menés au seuil usuel de 5~\%, le seuil de Bonferroni restant
réservé au verdict final.

% --- Le pré-test : y a-t-il quelque chose à corriger ? ----------------
\subsubsection*{Test d'indépendance de Box et Pierce}

Avant de corriger, encore faut-il qu'il y ait quelque chose à corriger. Le
raisonnement du recouvrement rend l'autocorrélation plausible, il ne la rend
pas certaine, et elle peut différer d'une stratégie à l'autre. Le protocole
commence donc par poser directement la question à chaque série, celle de
savoir si les mois sont indépendants. Le test de Box et
Pierce~\cite{boxpierce1970} y répond sur la série brute. Sa statistique cumule
les six premières autocorrélations, mises au carré pour que les positives et
les négatives ne se compensent pas, puis rapportées à la longueur de la
série,
\[
  S_{\mathrm{BP}} \;=\; T \sum_{h=1}^{6} \hat\rho(h)^2 ,
\]
où \(T\) est le nombre de mois de la série et \(\hat\rho(h)\)
l'autocorrélation d'ordre \(h\), c'est-à-dire la corrélation entre la série
et elle-même décalée de \(h\) mois. On oppose alors~:
\[
  \begin{cases}
    \Hzero \;(S_{\mathrm{BP}} \text{ petit}) : \text{les mois sont indépendants,}\\[2pt]
    \Hun \;(S_{\mathrm{BP}} \text{ grand}) : \text{au moins un décalage porte une autocorrélation.}
  \end{cases}
\]
C'est \Hzero{} qui porte ici la propriété recherchée.

Sous \Hzero{}, chaque $\hat\rho(h)$ n'est que du bruit d'échantillonnage
d'ordre $1/\sqrt{T}$, et la somme $S_{\mathrm{BP}}$ suit alors une loi du
khi-deux à six degrés de liberté. On calcule donc la p-value de notre
$S_{\mathrm{BP}}$ observé, et si elle est trop faible, l'indépendance cesse
d'être une explication crédible.

Si l'indépendance n'est pas rejetée, la condition du test par agrégation
mensuelle est remplie telle quelle. Son verdict tient et aucune correction
n'est requise, car on ne corrige que ce qui a été mesuré, exactement comme
pour l'effet de sondage. Si elle est rejetée, l'autocorrélation est réelle et
la suite de la méthode s'applique.

Cette correction suppose que tous les mois soient pareillement exposés à ceux
qui les précèdent. Or cette affirmation n'est vraie que si la série est
stationnaire, c'est-à-dire si ses caractéristiques statistiques ne changent
pas au fil du temps. Deux tests seront donc menés pour vérifier la
stationnarité de la série de chaque stratégie concernée.

Si une série ne l'était pas, on sortirait du domaine de validité de la
méthode. Le protocole la signale et s'abstient de conclure plutôt que de
produire un verdict fondé sur un impact moyen qui n'aurait pas de sens. Aucune
de nos stratégies n'est dans ce cas, mais la situation peut se présenter pour
une stratégie à venir.

% --- Stationnarité, versant 1 -----------------------------------------
\subsubsection*{Test de stationnarité de Dickey et Fuller augmenté}

Ce test, comme le suivant, se calcule directement sur la série. Il consiste à
estimer, par une simple régression du mois courant sur le mois précédent, un
coefficient que nous noterons \(\phi\), selon
\(x_t = \phi\, x_{t-1} + \eta_t\).
En remplaçant \(x_{t-1}\) par sa propre expression, puis en recommençant, on
voit ce que devient un choc au fil du temps,
\[
  x_t = \eta_t + \phi\,\eta_{t-1} + \phi^2\,\eta_{t-2} + \phi^3\,\eta_{t-3} + \cdots
\]
Le mois courant est donc la somme de tous les chocs passés, chacun pondéré par
\(\phi\) élevé à son ancienneté. Si \(\phi\) est inférieur à 1, ces puissances
tendent vers zéro, l'impact d'un mois sur ceux qui suivent décroît, et la série
oublie ses chocs pour revenir vers sa moyenne, ce qui la rend stationnaire. Si
\(\phi\) vaut 1, toutes les puissances valent 1, chaque choc s'ajoute
définitivement au niveau, et la série dérive sans repère, donc sans être
stationnaire. On oppose ainsi

\[
  \begin{cases}
    \Hzero \;(\phi = 1) : \text{les chocs s'accumulent, la série n'est pas stationnaire,}\\[2pt]
    \Hun \;(\phi < 1) : \text{les chocs s'estompent, la série est stationnaire.}
  \end{cases}
\]

Le test est unilatéral, car un \(\phi\) supérieur à 1 décrirait une série
explosive, cas sans objet ici. C'est donc le rejet qui conclut à la
stationnarité.

Comme pour le test de Student, la statistique compte le nombre d'erreurs-types
qui séparent le \(\hat\phi\) estimé de la valeur 1,

\[
  t_{\mathrm{DF}} = \frac{\hat\phi - 1}{\widehat{\mathrm{se}}(\hat\phi)},
\]


où \(\hat\phi\) est le coefficient estimé par la régression et
\(\widehat{\mathrm{se}}(\hat\phi)\) son erreur-type, c'est-à-dire
l'incertitude qui entoure cette estimation. Plus \(t_{\mathrm{DF}}\) est
négatif, plus \(\hat\phi\) s'écarte de 1 vers le bas et plus \Hzero{} devient
implausible.

Cette statistique, que les auteurs eux-mêmes nomment rapport de Student, n'en
suit pourtant pas la loi~\cite{dickeyfuller1979}. La raison tient à \Hzero{},
sous laquelle la série n'est justement pas stationnaire, sa variance croît avec
le temps, et les conditions qui donnent la loi de Student ne sont plus réunies.
Sa loi est alors celle que Dickey et Fuller ont établie et tabulée, et qui
porte leur nom. La lecture reste ensuite la même, la p-value étant confrontée
au seuil et le rejet prononcé si elle est trop faible.

Ce test suppose que les résidus de la régression ne soient plus autocorrélés.
On emploie pour cela sa version dite augmentée, qui ajoute des mois passés
supplémentaires dans la régression jusqu'à les nettoyer. Reste à savoir
combien en retenir, car trop peu laisserait la mémoire non absorbée, tandis
que chaque coefficient supplémentaire, estimé sur les mêmes données, ajoute du
bruit. C'est le critère d'information d'Akaike~\cite{akaike1974} qui arbitre,
\[
  \mathrm{AIC}_k = 2\ln \hat\sigma_k + \frac{k}{T},
\]
où \(\hat\sigma_k\) mesure l'erreur résiduelle du modèle à \(k\) retards. Le
premier terme diminue quand le modèle colle mieux aux données, tandis que le
second augmente avec le nombre de coefficients. On retient le \(k\) qui
minimise la somme, de sorte qu'un retard n'est conservé que s'il apporte plus
qu'il ne coûte. Sur nos séries, ce nombre va de zéro à neuf selon les
stratégies, et la condition se trouve ainsi prise en charge par le test
lui-même.

% --- Stationnarité, versant 2 -----------------------------------------
\subsubsection*{Test de stationnarité KPSS}

À l'inverse, le test de Kwiatkowski et ses coauteurs, désigné dans la
littérature par le sigle KPSS~\cite{kpss1992}, cumule les écarts à la moyenne
au fil du temps. Si la série oscille autour d'un niveau fixe, ces cumuls se
compensent et restent bornés, tandis qu'ils s'emballent dès que la série
dérive. Ses hypothèses sont inversées par rapport au test précédent,
\[
  \begin{cases}
    \Hzero \;(\text{cumuls bornés}) : \text{la série est stationnaire,}\\[2pt]
    \Hun \;(\text{cumuls qui s'emballent}) : \text{la série n'est pas stationnaire.}
  \end{cases}
\]
Comme pour le test de normalité vu plus haut, c'est ici \Hzero{} qui porte la
propriété recherchée, si bien que c'est l'absence de rejet qui conclut à la
stationnarité.

Sa statistique somme les carrés de ces cumuls, rapportés à la longueur de la
série et à sa variance de long terme, de sorte qu'elle reste petite quand les
cumuls sont bornés et explose quand ils dérivent. Sa loi sous \Hzero{} est là
aussi tabulée par les auteurs. On rejette donc la stationnarité quand la
statistique dépasse la valeur critique, c'est-à-dire quand la p-value passe
sous le seuil.

Ce test demande enfin de préciser autour de quoi la stationnarité est jugée.
Nous la testons autour d'une constante et non d'une tendance, car rien ne
laisse attendre qu'un avantage dérive régulièrement au fil des années.

% --- La règle de décision sur la stationnarité ------------------------
\subsubsection*{Concordance des deux tests}

Les deux hypothèses nulles étant opposées, chacun des tests protège l'autre de
sa propre faiblesse. Ne pas rejeter n'a en effet jamais valeur de preuve, et un
test seul ne distingue pas une série réellement stationnaire d'une série sur
laquelle il manque simplement de puissance. Nous ne conclurons donc que
lorsque les deux concordent, et s'ils divergent, le diagnostic sera tenu pour
non concluant plutôt que pour défavorable.

Ces deux tests ne demandent aucun modèle, et il faut bien l'entendre ainsi. Le
\(\phi\) de Dickey et Fuller ne sert qu'à trancher sur la racine unitaire et
n'est pas conservé. Les impacts que nous recherchons sont d'une autre nature,
et la stationnarité, maintenant acquise, nous donne le droit de les estimer.

% --- L'ordre du modèle : combien de mois comptent ---------------------
\subsubsection*{Autocorrélation partielle et choix de l'ordre}

On mesure l'impact propre de chaque décalage \(h\) par l'autocorrélation
partielle \(\hat\tau(h)\). Reste à décider quels décalages comptent vraiment.
On procède décalage par décalage, en opposant
\[
  \begin{cases}
    \Hzero \;\bigl(\tau(h) = 0\bigr) : \text{le mois } t-h \text{ n'a pas d'impact propre,}\\[2pt]
    \Hun \;\bigl(\tau(h) \neq 0\bigr) : \text{il en a un.}
  \end{cases}
\]
Sous \Hzero{}, \(\hat\tau(h)\) n'est que du bruit d'échantillonnage,
approximativement normal et d'écart-type \(1/\sqrt{T}\), si bien que la
statistique \(\sqrt{T}\,\hat\tau(h)\) suit une loi normale centrée réduite. On
en tire, comme pour les tests précédents, la p-value de notre \(\hat\tau(h)\)
observé, et si elle passe sous le seuil de 5~\%, l'absence d'impact cesse
d'être une explication crédible et le décalage est retenu. Cela revient à
comparer \(\bigl|\hat\tau(h)\bigr|\) à \(q_{0{,}975}/\sqrt{T}\), avec
\(q_{0{,}975} \approx 1{,}96\), ce qui définit la bande de significativité que
l'on trace sur les corrélogrammes.

L'ordre retenu \(\hat p\) est le dernier décalage significatif. Si des
décalages intermédiaires ne le sont pas, ils sont conservés malgré tout, car
on préfère surestimer la mémoire plutôt que la tronquer.

La recherche est plafonnée à six mois. Ce plafond tient au mécanisme même de
la dépendance, qui vient du recouvrement des périodes de détention et ne peut
donc porter plus loin que la durée des positions, lesquelles dépassent
rarement six mois. Chercher au-delà reviendrait à chercher un effet sans
cause.

Cet argument reste toutefois une plausibilité et non une preuve. Nous n'en
resterons donc pas là, puisque la validation du modèle, présentée plus bas,
vérifiera qu'aucune mémoire ne subsiste au-delà de ce plafond.

% --- L'estimation des poids -------------------------------------------
\subsubsection*{Équations de Yule-Walker}

L'ordre étant connu, on écrit le modèle autorégressif sur la série centrée, le
mois courant s'exprimant comme une somme pondérée des \(\hat p\) mois
précédents, augmentée d'un bruit,
\[
  x_t = \Phi_1\, x_{t-1} + \cdots + \Phi_{\hat p}\, x_{t-\hat p} + \eta_t ,
\]
où \(x_t\) est l'avantage moyen du mois \(t\) (écarté de la moyenne de la
série), \(\Phi_i\) le poids du mois \(t-i\) — l'impact propre que nous
cherchions —, et \(\eta_t\) le bruit, c'est-à-dire la part du mois courant
qu'aucun mois passé n'explique.
Ces poids ne se lisent pas directement dans les données. Les autocorrélations,
en revanche, se mesurent sans difficulté, puisqu'il suffit de comparer la série
à elle-même décalée. Les équations de Yule et Walker font le lien entre les
deux. On les obtient en multipliant le modèle par \(x_{t-h}\), le mois situé
\(h\) crans en arrière, puis en prenant l'espérance. Le bruit \(\eta_t\)
disparaît alors, étant indépendant de tout ce qui précède, et il reste, pour
chaque \(h = 1, \ldots, \hat p\),
\[
  \rho(h) = \Phi_1\,\rho(h-1) + \Phi_2\,\rho(h-2) + \cdots + \Phi_{\hat p}\,\rho(h-\hat p),
\]
avec \(\rho(0) = 1\) et \(\rho(-k) = \rho(k)\), une série étant liée à
elle-même de la même façon dans les deux sens. Chaque autocorrélation
s'exprime ainsi en fonction des poids recherchés. Comme il y a autant
d'équations que de poids inconnus, on y reporte les autocorrélations mesurées,
on résout, et l'on remonte aux poids. Le cas d'un seul mois de mémoire est
éclairant, car le système se réduit alors à \(\rho(1) = \Phi_1\), le poids
étant exactement l'autocorrélation mesurée.

Une équation supplémentaire, celle du décalage nul, donne l'amplitude du
bruit,
\[
  \sigma^2_\eta = \gamma(0)\Bigl(1 - \sum_{i=1}^{\hat p} \Phi_i\,\rho(i)\Bigr),
\]
soit la variance de la série moins la part qu'expliquent les mois passés,
autrement dit ce qui reste imprévisible.

Ce système n'a de sens que si la série est stationnaire. Parler de
l'autocorrélation à deux mois d'écart suppose en effet qu'il en existe une
seule valeur, la même en 2010 qu'en 2020. Sans cette garantie, il en faudrait
une par époque, et le système n'aurait plus de solution unique.

C'est la somme de ces poids qui nous intéresse, car elle mesure la persistance
totale de la série et dira de combien l'erreur-type du test mensuel doit être
corrigée.

% --- La correction et le verdict --------------------------------------
\subsubsection*{Variance de long terme et test de Student corrigé}

Le modèle fournit alors une variance de la moyenne qui tient compte de la
persistance, appelée variance de long terme,
\[
  \gamma_{\mathrm{lr}}
     = \frac{\sigma^2_\eta}{\bigl(1-\sum_{i=1}^{\hat p}\Phi_i\bigr)^{2}},
  \qquad
  \widehat{\mathrm{Var}}(\bar x) \approx \frac{\gamma_{\mathrm{lr}}}{T},
  \qquad
  \mathrm{DEFF}_t = \frac{\gamma_{\mathrm{lr}}}{\gamma(0)},
\]
où \(\sigma^2_\eta\) et les \(\Phi_i\) sont le bruit et les poids estimés
ci-dessus, \(T\) le nombre de mois et \(\gamma(0)\) la variance de la série.
Cette variance de long terme additionne toutes les autocovariances, là où la
formule usuelle en \(s^2/T\) suppose au contraire des mois sans lien. Si les
poids sont nuls, le dénominateur vaut 1 et l'on retrouve la formule usuelle,
tandis qu'une persistance forte rend la correction d'autant plus importante.

Le rapport \(\mathrm{DEFF}_t\) est l'exact analogue temporel de l'effet de
sondage de la première méthode. C'est un facteur qui dit combien de fois la
variance de la moyenne est sous-estimée lorsqu'on ignore la dépendance, cette
fois sur l'axe du temps, et obtenu par un modèle puisque la dépendance y est
ordonnée.

Le test de Student de la porte de la période est alors repris tel quel, avec
deux changements. L'erreur-type devient \(\sqrt{\gamma_{\mathrm{lr}}/T}\), et
le nombre de degrés de liberté est ramené à la taille effective
\(T_{\mathrm{eff}} = T/\mathrm{DEFF}_t\), comme pour l'effet de sondage. La
statistique diminue mécaniquement, si bien qu'une stratégie peut
franchir la porte de la période et échouer ici. Le verdict se lit au même
seuil de Bonferroni que les étapes précédentes.

% --- La validation finale ---------------------------------------------
\subsubsection*{Validation du modèle~: Box et Pierce sur les résidus}

Reste à vérifier que le modèle a bien capté toute l'autocorrélation, sans quoi
la correction serait fondée sur une mémoire mal décrite. On dispose pour cela
des résidus \(\hat\eta_t\), c'est-à-dire de ce que le modèle n'explique pas,
soit pour chaque mois la valeur observée diminuée de celle que les mois
précédents laissaient attendre. Le raisonnement est alors le même qu'au début
de la méthode, mais appliqué à ce qui reste. Si le modèle a bien absorbé toute
la mémoire, les résidus n'en contiennent plus et sont indépendants d'un mois à
l'autre. On leur pose donc exactement la question posée à la série brute, avec
le même test,
\[
  \begin{cases}
    \Hzero \;\bigl(S_{\mathrm{BP}} \text{ petit}\bigr) : \text{les résidus sont indépendants, le modèle a tout absorbé,}\\[2pt]
    \Hun \;\bigl(S_{\mathrm{BP}} \text{ grand}\bigr) : \text{il subsiste de la mémoire, le modèle est inadéquat.}
  \end{cases}
\]
La démarche est donc circulaire au bon sens du terme, le premier test
demandant s'il y a de la mémoire et celui-ci s'il en reste. Deux différences
seulement les séparent. L'inspection porte sur quinze décalages et non plus
six, comme le recommande la pratique usuelle~\cite{boxpierce1970} qui situe ce
nombre entre quinze et vingt, de sorte que le plafond de six mois n'a rien pu
laisser échapper. Par ailleurs, comme \(\hat p\) coefficients ont été estimés
sur ces mêmes données, ce qui rend les résidus artificiellement un peu plus
propres, on retranche le nombre de paramètres du modèle, et la statistique
suit sous \Hzero{} une loi du khi-deux à \(15 - \hat p\) degrés de liberté.

Ne pas rejeter valide donc la correction. Un rejet signalerait au contraire un
ordre trop faible ou une forme inadaptée, et le protocole s'abstiendrait alors
de conclure, comme pour la stationnarité.

Au terme de cette méthode, la condition d'indépendance entre mois, que la
porte de la période supposait sans la vérifier, a donc été soit établie, soit
prise en charge par une correction elle-même contrôlée. Rappelons que les
conditions ne s'appliquent qu'à la branche qui les appelle. Une stratégie dont
l'indépendance mensuelle n'est pas rejetée conserve le verdict de la porte de
la période sans qu'on ait à exiger d'elle ni stationnarité ni modèle,
puisqu'aucune correction ne lui est appliquée.

\subsection{Résultats}

Nous avons vu précédemment que seules trois stratégies franchissaient le
premier filtre, celui du test de Student naïf qui supposait l'indépendance.
Examinons maintenant les résultats du protocole sur ces stratégies.

% --- 1b : diagnostic + correction DEFF (les deux « survivantes ») -----
\begin{table}[H]
\centering
\caption{Diagnostic de dépendance intra-titre et test corrigé (étape 1b).}
\label{tab:b1-dep-deff}
\small
\begin{tabular}{@{}lrrrccc@{}}
\toprule
Stratégie & $\hat\rho$ & DEFF & $n_{\text{eff}}$ & $p$ naïve & $p$ corrigée & $p$ bootstrap \\
\midrule
\texttt{oracle} (témoin) & $0{,}053$ & $1{,}67$ & $4\,076$ & $\approx 0$ & $\approx 0$ & $<2{,}5\times10^{-4}$ \\
\midrule
\texttt{rsi\_classic} & $0{,}000$ & $1{,}00$ & $7\,117$ & $1{,}2\times10^{-8}$ & $1{,}2\times10^{-8}$ & $<2{,}5\times10^{-4}$ \\
\texttt{rsi\_strict}  & $0{,}134$ & $1{,}15$ & $848$    & $5{,}5\times10^{-10}$ & $7\times10^{-9}$     & $<2{,}5\times10^{-4}$ \\
\multicolumn{7}{c}{\footnotesize (autres stratégies~: verdict « non » inchangé, $\hat\rho \le 0{,}09$)} \\
\bottomrule
\end{tabular}
\end{table}
Le $\hat\rho$ intra-titre mesuré au chapitre~\ref{chap:donnees} était quasi nul
pour l'ensemble des stratégies, à l'exception de \texttt{rsi\_strict} avec
$0{,}134$. La correction par l'effet de sondage, qui gonfle la variance en
proportion de cette dépendance, ne change donc presque rien.
\texttt{rsi\_classic} et \texttt{rsi\_strict} franchissent toutes deux cette
première étape, et le rééchantillonnage par grappes, qui ne fait aucune
hypothèse sur la forme de la dépendance, confirme ce verdict.

Le protocole se poursuit donc sur ces trois stratégies.

% --- 1c : les deux portes (TABLEAU CENTRAL de la section) --------------
\begin{table}[H]
\centering
\caption{Le test des deux portes~: une unité réelle, un vote (étape 1c).}
\label{tab:b1-dep-portes}
\small
\begin{tabular}{@{}lrrrrc@{}}
\toprule
 & \multicolumn{2}{c}{Porte action ($k$ titres)} & \multicolumn{2}{c}{Porte période ($k$ mois)} & \\
\cmidrule(lr){2-3}\cmidrule(lr){4-5}
Stratégie & $k$ & $p_A$ & $k$ & $p_B$ & Verdict \\
\midrule
\texttt{oracle} (témoin) & 501 & $\approx 0$          & 193 & $\approx 0$ & \textbf{OUI} \\
\midrule
\texttt{rsi\_classic}  & 501 & $0{,}059$            & 195 & $0{,}039$ & non \\
\texttt{rsi\_strict}   & 456 & $3{,}9\times10^{-4}$ & 138 & $0{,}26$  & non \\
\texttt{rsi\_trend}    & 495 & $0{,}22$             & 175 & $0{,}29$  & non \\
\texttt{sr\_breakdown} & 457 & $0{,}37$             & 184 & $0{,}94$  & non \\
\multicolumn{6}{c}{\footnotesize (six autres stratégies~: $p_A \ge 0{,}74$, toutes « non »)} \\
\bottomrule
\end{tabular}
\end{table}
Ce tableau change la lecture du chapitre. L'oracle franchit les deux portes de
très loin, ses deux p-values étant pratiquement nulles. Son avantage, construit
pour être réel, se retrouve aussi bien titre par titre que mois par mois.
C'est le résultat attendu, mais il n'en est pas moins important, car il montre
que l'outil ne rejette pas tout par excès de prudence et qu'il sait reconnaître
un avantage lorsqu'il est effectivement présent.

\texttt{rsi\_classic}, en revanche, tombe aux deux portes, avec
$p_A = 0{,}059$ et $p_B = 0{,}039$, loin du seuil de Bonferroni fixé à
$0{,}0045$. Les 7\,117 opérations de cette stratégie ne constituaient pas
7\,117 témoignages indépendants, mais un nombre bien plus restreint de titres
et de mois, chacun compté plusieurs fois. Dès que chaque titre, puis chaque
mois, ne dispose plus que d'une seule voix, l'avantage s'efface.

\texttt{rsi\_strict} présente un comportement différent. Elle franchit
brillamment la porte de l'action, avec $p_A = 3{,}9\times10^{-4}$, ce qui
montre que son avantage est bien partagé par un grand nombre de titres, mais
elle échoue à la porte de la période avec $p_B = 0{,}26$. Son avantage est
donc réel titre par titre, mais épisodique dans le temps, sans régularité
démontrable d'un mois à l'autre. En l'état, il n'est pas exploitable.

Le protocole écarte ces deux stratégies, tout en signalant qu'elles présentent
des résultats remarquables sur certains titres ou certains mois. Peut-être
existe-t-il un facteur qui expliquerait ces résultats, mais l'établir
demanderait une étude spécifique, qui devrait se garder du sur-apprentissage.
En l'état, on ne peut pas les laisser passer.

En dehors du témoin, plus aucune stratégie ne franchit les deux portes. Cela
n'a rien d'étonnant, puisque nous opérons sur les cinq cents plus grandes
capitalisations américaines, dont la forte tendance à croître exige un
résultat extrême pour dépasser le seuil de l'improbable.



% --- 1d : la porte « mois » corrigée de l'autocorrélation -------------
\begin{table}[H]
\centering
\caption{Série mensuelle des edges~: autocorrélation et test corrigé (étape 1d).}
\label{tab:b1-dep-mensuel}
\small
\begin{tabular}{@{}lrrrrrc@{}}
\toprule
Stratégie & $T$ & $\hat\rho(1)$ & $\hat p$ (AR) & $\mathrm{DEFF}_t$ & $p$ corrigée & $p$ Box--Pierce \\
\midrule
\texttt{oracle} (témoin) & 193 & $0{,}18$ & 0 & $1{,}0$ & $9{,}4\times10^{-85}$ & $0{,}16$ \\
\midrule
\texttt{ma\_crossover} & 188 & $0{,}56$ & 3 & $6{,}4$ & $0{,}94$ & $0{,}14$ \\
\texttt{dt\_top}       & 197 & $0{,}38$ & 1 & $2{,}2$ & $0{,}97$ & $0{,}69$ \\
\texttt{db\_bottom}    & 197 & $0{,}36$ & 1 & $2{,}1$ & $0{,}82$ & $0{,}90$ \\
\texttt{hs\_inverse}   & 186 & $0{,}31$ & 1 & $1{,}9$ & $0{,}91$ & $0{,}93$ \\
\texttt{rsi\_classic}  & 195 & $0{,}23$ & 1 & $1{,}6$ & $0{,}08$ & $0{,}93$ \\
\texttt{rsi\_strict}   & 138 & $0{,}03$ & 0 & $1{,}0$ & $0{,}26$ & $0{,}28$ \\
\multicolumn{7}{c}{\footnotesize (quatre autres stratégies~: $\mathrm{DEFF}_t \le 2{,}4$, aucune significative)} \\
\bottomrule
\end{tabular}
\end{table}
Cette dernière étape ne concerne que les stratégies encore en lice, celles qui
auraient franchi les deux portes. Pour chacune, on se demande d'abord si les
mois voisins sont réellement liés. Le test de Box et Pierce, appliqué à la
série brute, tranche entre une série sans mémoire, qui conserve alors le
verdict de la porte de la période sans correction, et une série autocorrélée,
qui en appelle une.

Le tableau montre que la moitié des séries seulement sont autocorrélées et
suivent la branche corrigée, l'autocorrélation d'ordre un atteignant $0{,}56$
pour \texttt{ma\_crossover}. Les autres, dont l'oracle et \texttt{rsi\_strict},
ne rejettent pas l'indépendance des mois, et leur verdict tient donc tel quel.

Pour les séries autocorrélées, la correction n'est légitime que sous
stationnarité. Le croisement des deux tests l'assure ici pour toutes, le test
de Dickey et Fuller rejetant la racine unitaire avec $p \le 0{,}005$, tandis
que le KPSS ne rejette pas la stationnarité. Après correction, aucun verdict
ne bascule. La variance mensuelle était bien sous-estimée, d'un facteur allant
jusqu'à $6{,}4$ pour \texttt{ma\_crossover}, mais aucune stratégie ne devient
significative pour autant. Comme attendu, l'oracle franchit lui aussi cette
dernière étape.

Une réserve doit être signalée. Les mois durant lesquels une stratégie n'a
produit aucune opération sont traités comme consécutifs, ce qui revient à
ignorer les trous de la série. La couverture dépasse 88~\% pour presque toutes
les stratégies, mais tombe à 71~\% pour \texttt{rsi\_strict}, dont la série
est donc la plus concernée par cette approximation.



\subsection{Positionnement dans la littérature financière}
\label{sec:b1-positionnement}
Les outils de ce chapitre ont été construits à partir des seuls concepts vus
en cours, qu'il s'agisse de la théorie des sondages, de l'analyse de la
variance ou des séries chronologiques. Or ils coïncident avec les standards de
l'économétrie financière, ce qui indique que la démarche rejoint d'elle-même
la pratique établie.

L'agrégation par mois, qui calcule une statistique par période avant de mener
l'inférence sur la suite de ces périodes, est exactement la démarche de Fama
et MacBeth~\cite{famamacbeth1973}, devenue un standard du domaine.

La variance de long terme estimée à partir du modèle autorégressif poursuit
quant à elle le même objectif que les erreurs-types robustes à
l'autocorrélation de Newey et West~\cite{neweywest1987}, dont elle constitue
une version paramétrique, validée ici par le test de Box et Pierce.

AUCUNE des onze stratégies, hormis le témoin, ne
démontre d'avantage réel. Le verdict de l'étape~1 ne tenait que par
une condition non vérifiée. L'outil a donc écarté toutes les stratégies
dont il ne peut affirmer l'avantage sans que cela puisse relever de la chance.
% ------------------------------------------------------------------
% TODO (mécanique, hors rédaction — à faire par le chat sur demande) :


% =====================================================================
\section{Synthèse du chapitre}
\label{sec:b1-synthese}
Ce chapitre a procédé par une escalade de prudence. Nous sommes partis du test
de Student, le point de départ le plus simple pour une telle analyse. Malgré
ses hypothèses simplificatrices, il suffisait déjà à écarter huit des onze
stratégies étudiées.

Il a ensuite fallu prendre en compte la dépendance, restée jusque-là hors du
champ du test. En reprenant les résultats du chapitre~\ref{chap:donnees}, qui
avait mis en évidence trois canaux se ramenant à deux portes, la dépendance
entre les opérations d'un même titre d'une part, celle liée au moment d'autre
part, nous avons accepté de perdre beaucoup d'information en agrégeant par
titre puis par mois.

Cette agrégation aurait suffi s'il n'était pas resté, en outre, la dépendance
entre mois voisins. Une dernière correction s'est donc imposée, qui a
elle-même exigé de vérifier la stationnarité des séries.

Deux résultats méritent d'être retenus. \texttt{rsi\_classic} apparaissait
gagnante au test initial, mais son avantage s'est révélé n'être qu'un artefact
de comptage, ses milliers d'opérations ne portant qu'un nombre bien plus
restreint d'informations réellement distinctes. \texttt{rsi\_strict}, elle,
possède un avantage authentiquement partagé entre les titres, mais trop
irrégulier dans le temps pour être retenu. Au terme de l'escalade, plus aucune
stratégie ne subsiste en dehors du témoin.

C'est précisément là que se joue la démonstration. Un outil de validation ne
se juge pas à sa capacité de produire des verdicts favorables, mais à celle de
démasquer ses propres erreurs tout en reconnaissant ce qui est vrai. Le témoin
a franchi chaque étage, les stratégies réelles y ont échoué l'une après
l'autre, et le faux positif du test initial a été mis au jour par la
vérification de ses propres conditions. Un verdict négatif rendu proprement
vaut ici démonstration que l'outil fonctionne.

Le chapitre suivant expose les limites de ce protocole, qui pourra
certainement être perfectionné avec le temps.